\documentclass[twocolumn]{aastex631}
\begin{document}
\title{A residual reionization ionizing-photon-budget shortfall for star-forming galaxies at z\~{}9 (o32 systematic corner)}
\author{NebulaMind Lab (autonomous pipeline)}
\affiliation{NebulaMind Lab, \url{https://lab.nebulamind.net}}
\begin{abstract}
Reionization ionizing-photon-budget reconciliation at z\~{}9: star-forming galaxies require f\_esc=0.390 (+0.393/-0.200) to close the budget (Madau-Dickinson SFRD, log xi\_ion=25.5+/-0.15, clumping C=2-5, JWST-SFRD tail), versus indirect-proxy-inferred f\_esc=0.080 (+0.146/-0.051) from LzLCS O32/beta calibrations. Median delta(required-inferred)=+0.280 dex-frac (16-84\%: +0.055 to +0.673); 89\% of the systematic MC shows a shortfall. Conclusion: a genuine SHORTFALL remains, and the sign holds under both O32 and beta calibrations. The result is bounded by the xi\_ion x clumping x proxy-calibration systematic, not by statistics. Generated autonomously from public data (jwst) via the NebulaMind Lab runner: a bounded, reproducible, descriptive study.
\end{abstract}
\section{Introduction}
Recent studies have highlighted concerns about a potential crisis in our understanding of reionization, suggesting that current models may not account for the necessary ionizing photons to achieve reionization [Muñoz2024]. This issue is further complicated by the increased demands on ionizing sources due to absorption-dominated reionization [Davies2021]. To address this, we revisit the ionizing-photon-budget using a literature-anchored budget calculation. Our approach builds upon previous work that aimed to conserve ionizing photons during reionization [Park2022] and assesses the galaxy ionizing photon budget at high redshifts [Duncan2015].
\section{Data and method}
In our analysis, we rely on established values from the literature: the cosmic star formation rate density (SFRD) is based on the Madau \& Dickinson (2014) analytic fitting function. We adopt published calibrations for xi\_ion and O32/beta f\_esc proxy from LzLCS [Chisholm+22, Flury+22; Simmonds+24]. Notably, we do not utilize survey catalog data or observations from JWST, SDSS, or TNG in this study. Instead, our method focuses on reconciling systematics across published literature values using an ionizing-photon-budget approach.
\section{Result}
Our reconciliation of the reionization ionizing-photon-budget at z\~{}9 suggests that star-forming galaxies may require a higher escape fraction (f\_esc=0.390 +0.393/-0.200) to close the budget, considering the Madau-Dickinson SFRD, log xi\_ion=25.5±0.15, and clumping factor C=2-5. However, it is crucial to acknowledge that this result depends on assumptions about these parameters, which have associated uncertainties that may affect its validity. In contrast, indirect-proxy-inferred f\_esc from LzLCS O32/beta calibrations is 0.080 (+0.146/-0.051). The median difference between the required and inferred escape fractions is +0.280 dex-frac (16-84\%: +0.055 to +0.673), with 89\% of systematic Monte Carlo simulations showing a shortfall.
\begin{figure}\centering\includegraphics[width=\columnwidth]{result.png}\caption{A residual reionization ionizing-photon-budget shortfall for star-forming galaxies at z\~{}9 (o32 systematic corner)}\end{figure}
\section{Caveats}
It is essential to acknowledge the limitations of our approach. Our analysis relies on automated, single-selection, uncalibrated measurements, which may introduce biases and uncertainties. The result is bounded by systematics related to xi\_ion, clumping factor, and proxy-calibration, rather than statistical errors. Additionally, our study does not incorporate new observational data from JWST or other sources to validate the assumptions made about SFRD and escape fractions at high redshifts. Therefore, while our findings suggest a potential shortfall in ionizing photons, further research is needed to refine these estimates, address the underlying systematics, and incorporate new observational data to strengthen the analysis. A more comprehensive understanding of the reionization-photon-budget crisis will require additional studies that directly measure key parameters at high redshifts.
\section*{References}
\footnotesize
[Muoz2024] Muñoz et al. (2024). Reionization after JWST: a photon budget crisis?. bibcode 2024MNRAS.535L..37M \par [Davies2021] Davies et al. (2021). The Predicament of Absorption-dominated Reionization: Increased Demands on Ionizing Sources. bibcode 2021ApJ...918L..35D \par [Park2022] Park et al. (2022). Calibrating excursion set reionization models to approximately conserve ionizing photons. bibcode 2022MNRAS.517..192P \par [Duncan2015] Duncan et al. (2015). Powering reionization: assessing the galaxy ionizing photon budget at z < 10. bibcode 2015MNRAS.451.2030D \par [Madau2017] Madau (2017). Cosmic Reionization after Planck and before JWST: An Analytic Approach. bibcode 2017ApJ...851...50M

\end{document}
